% Generado por generar_capitulos.py a partir de ../.tex8
\chapter{Planteamiento del Problema}\label{cap:pp}
\section{Antecedentes}
% <- 01 sin la pregunta (los primeros párrafos se repiten en la Introducción)
Cuando un grupo de personas trabaja de manera coordinada durante el tiempo suficiente,
puede recordar más de lo que cualquiera de sus miembros recordaría por separado.
Wegner \cite{wegner1987transactive} llamó a este fenómeno
\textbf{memoria transactiva}: una forma de memoria compartida en la que cada individuo
almacena no solo sus propios conocimientos, sino también un \emph{directorio} que indica quién es reconocido
dentro del grupo como referente en cada dominio.
Cuando un miembro necesita información que no conoce, consulta ese directorio,
determina a quién preguntar y delega la recuperación.
El grupo funciona gracias a esa especialización: el conocimiento no necesita repetirse
en cada miembro, sino que se reparte entre ellos.

El modelo de Wegner identifica tres procesos principales \cite{wegner1995computer}.
Primero, la \emph{actualización del directorio}: tras cada interacción, los
miembros que la presenciaron registran quién supo responder.
Segundo, la \emph{asignación de información}: la información nueva se envía
hacia el miembro reconocido como referencia en ese dominio.
Tercero, la \emph{coordinación de recuperación}: las consultas futuras viajan directamente
a ese referente sin pasar por un intermediario.
En este trabajo, la transición de una \emph{fase temprana} mediada a una \emph{fase
madura} punto a punto es la señal de que el sistema transactivo ya quedó formado.

Las memorias asociativas entrópicas de Pineda y Morales (Sección~\ref{sec:eam}) guardan
lo que registran como conteos de frecuencia: aceptan una entrada cuando sus rasgos
tienen soporte en lo registrado, recuerdan muestreando esos conteos y rechazan lo que no
contienen, sin ningún clasificador entrenado.

El sistema construido, \textbf{EAM-TMS} (\textit{Entropic Associative Memory --
Transactive Memory System}: sistema de memoria transactiva sobre memorias asociativas
entrópicas), se compone de ocho módulos especialistas, uno por cada categoría del
conjunto ETH-80 \cite{leibe2003eth80}: \texttt{apple}, \texttt{car}, \texttt{cow},
\texttt{cup}, \texttt{dog}, \texttt{horse}, \texttt{pear} y \texttt{tomato}.
Cada módulo guarda las palabras y las imágenes de su categoría en una sola memoria de
contenido y lleva dos directorios, uno textual y uno visual, donde anota las
transacciones en que participa.
Un mediador central denominado TME (\textit{Transactive Memory Expert}) coordina la fase
temprana y registra en sus propios directorios qué módulo responde a cada tipo de
consulta.
Conforme se llenan los directorios de los módulos, el TME deja de hacer falta: en la fase
madura los módulos se reenvían entre sí las consultas, sean frases o imágenes, sin
intervención del mediador.

Ninguna decisión del sistema sale de un clasificador entrenado ni de una optimización
numérica: registrar, reconocer, recuperar y enrutar son operaciones definidas
formalmente sobre las Memorias Asociativas Entrópicas
de Pineda y Morales en su variante ponderada (W-EAM)
\cite{pineda2021entropic,pineda2022weighted} y su extensión heteroasociativa (EHAM)
\cite{morales2025missing}.
Las representaciones semánticas se obtienen con vectores de palabras fastText
\cite{bojanowski2017fasttext}, cuantizados \emph{por magnitud} en 16 niveles.
El vocabulario de llenado se extrae de ConceptNet 5.7 \cite{speer2017conceptnet}
y el texto se procesa con spaCy \cite{honnibal2020spacy}.
Las representaciones visuales las produce el codificador de un autocodificador ResNet18
\cite{he2016resnet}, inicializado con los pesos de ImageNet y entrenado sobre ETH-80
(Sección~\ref{sec:percepcion}); su salida se cuantiza y almacena.

\subsection*{Contribuciones principales}

Este trabajo aporta cinco contribuciones:

\begin{enumerate}
    \item \textbf{Un sistema transactivo de ocho especialistas.}
    Se construye un sistema de ocho módulos, uno por cada clase de ETH-80, cada uno con
    una sola memoria de contenido heteroasociativa y dos directorios, y se evalúa con un
    banco de 411 consultas ($\approx$50 por dominio).
    Los directorios, su lectura normalizada por número de registros (B1,
    Sección~\ref{sec:operaciones}), el rechazo de lo que ninguna memoria contiene y la
    coordinación punto a punto de la contribución~3 funcionan con los ocho dominios.

    \item \textbf{Cuantización por magnitud de las pistas semánticas.}
    Cada vector fastText se divide entre una escala global $S$, la misma en el llenado y
    en las consultas, se recorta a $[-1, 1]$ y se reparte en los 16 niveles
    (Sección~\ref{sec:pipeline-semantico}); una binarización por signo usaría solo 2.
    Con esta representación, la precisión temprana sobre el banco de 411 consultas es del
    90.0\% (Sección~\ref{sec:calibracion}).

    \item \textbf{Directorios perspectivales por módulo.}
    Cada módulo lleva dos directorios heteroasociativos, uno de texto y uno visual
    (Sección~\ref{sec:componentes}), y en ambos registra solo las transacciones en que
    participó, como módulo de entrada o como ganador: cada módulo sabe de los demás lo
    que presenció, y por eso los directorios se llaman \emph{perspectivales}.
    Una frase o una imagen puede entrar por \emph{cualquier} módulo, que suma su
    directorio y los de los módulos que conoce y la redirige punto a punto al
    especialista; es la coordinación de recuperación de Wegner, en las dos modalidades.

    \item \textbf{Auditoría de fidelidad.}
    Ninguna decisión proviene de mecanismos ajenos a las memorias: el vocabulario no
    contiene vectores sintéticos (de relleno, para etiquetas sin vector fastText), el
    código de Pineda y Morales permanece idéntico byte a byte (salvo una condición que
    evita cargar sus clasificadores y no toca ninguna operación de memoria), y todos los
    números reportados son trazables a corridas del sistema actual, sobre datos no
    fabricados.

    \item \textbf{Evaluación experimental.}
    Se miden nueve condiciones de lectura del directorio (con cuatro cantidades de
    consultas de formación y cinco semillas), el efecto de los parámetros nativos de la
    memoria ($\iota$ y $\kappa$), cuántas interacciones hacen falta para formar el
    directorio, cuánto contenido admite cada módulo, el comportamiento con imágenes, la
    comparación con una sola memoria para las ocho clases y la verificación de una sola
    memoria de contenido por módulo, junto con el costo computacional de reproducir
    todos los resultados.
\end{enumerate}

\subsection*{Estructura del documento}

La Sección~\ref{sec:marco_teorico} presenta los fundamentos: la memoria asociativa
entrópica de Pineda y Morales y su variante ponderada, su extensión heteroasociativa
(EHAM) y el modelo de memoria transactiva de Wegner.
La Sección~\ref{sec:arquitectura} describe la arquitectura EAM-TMS de ocho módulos, con
los diagramas de flujo del llenado, la fase temprana y la fase madura, y la
formalización de las operaciones del sistema.
La Sección~\ref{sec:metodologia} detalla la metodología experimental, la
Sección~\ref{sec:resultados} reporta los resultados, la Sección~\ref{sec:discusion} los
discute y señala las limitaciones, y la Sección~\ref{sec:conclusiones} concluye y señala
trabajo futuro.

\section{Definición del problema}
% <- 01, pregunta de investigación
Este trabajo pregunta si la organización que describe Wegner puede construirse solo con
estas memorias:

\begin{quote}
¿Puede un sistema de memoria transactiva construido sobre memorias asociativas
entrópicas repartir consultas de texto e imagen entre ocho especialistas, recordar el
contenido asociado y rechazar lo que no conoce, sin clasificadores supervisados? Y donde
aparezcan límites, ¿provienen del mecanismo asociativo o de las representaciones que lo
alimentan?
\end{quote}

\section{Objetivo general}\label{secc:og}
[PENDIENTE: redactar desde las contribuciones de 01]

\subsection{Objetivos específicos}
[PENDIENTE: redactar desde las contribuciones de 01]

\section{Justificación}
[PENDIENTE]

\section{Alcances y limitaciones}
[PENDIENTE: redactar alcances y limitaciones. Las limitaciones (a)--(j) están en la Sección~\ref{sec:limitaciones}.]
